% !PRINTED none
% Abschnitt-opening / part-title page: no page number in running head. Footer furniture
% "Weber, Algebra. I." and signature "29" are ignored per standing policy.

\begin{center}
\gesperrt{Thirteenth Section.}

\gesperrt{The Galois Theory.}
\end{center}

\begin{center}
\rule{3cm}{0.4pt}
\end{center}

\begin{center}
\S.\,139.

\gesperrt{The Concept of a Field.}
\end{center}

A system of numbers is called a \gesperrt{number field} when it is so complete and closed in itself that the four fundamental operations of reckoning (the four species) — \gesperrt{addition}, \gesperrt{subtraction}, \gesperrt{multiplication} and \gesperrt{division} — carried out with any numbers of the system, excepting division by zero, always lead to numbers that belong to the same system. This concept, which gives a classification of the kinds of numbers according to a natural point of view, was introduced by \gesperrt{Dedekind} (\gesperrt{Dirichlet-Dedekind}, Vorlesungen \"uber Zahlentheorie, 2nd ed. 1871, \S.\,159). It is of the greatest importance for algebra, and it is not a matter of indifference to have a distinctive and expressive name for it. The word \gesperrt{Zahlk\"orper} (number field) was formed by \gesperrt{Dedekind} on the model of numerous analogies, in which the word K\"orper (corpus, corps) in a similar way signifies a union of things belonging together, to which a certain completeness attaches.

The concept of the number field can be extended and carried over to all magnitudes with which one can reckon according to the rules of the four species; in particular, therefore, to rational functions of any variables.

We accordingly call a \gesperrt{function field} a system of functions of one or more variables of such a nature that in this system the reckonings with the four species can be carried out without limit and always lead to a definite function of the same system (always with the exception of division by zero). The variables
%CONT>>
